
@@M0@@
et l'égalité suivante dans @@M1@@ :
@@M2@@
En utilisant la filtration sur les @@M3@@ définie en 1.7, on trouve :
@@M4@@
Or @@M5@@~; on trouve donc :
@@M6@@
On a donc finalement exprimé @@M7@@ comme combinaison linéaire à coefficients dans @@M8@@ des puissances de @@M9@@.

1.9. Pour voir que @@M10@@ sont linéairement dépendants sur @@M11@@ on considère l'homomorphisme canonique surjectif : @@M12@@, dont le noyau @@M13@@ est localement libre de rang @@M14@@. On peut donc écrire :
@@M15@@
les opérations @@M16@@ sur @@M17@@ étant définies dans V 2.2. On obtient donc :
@@M18@@
Comme @@M19@@ est un faisceau inversible, @@M20@@, et par suite, @@M21@@, d'où finalement la relation :
@@M22@@

Montrons enfin que @@M23@@ sont indépendants sur @@M24@@\SGAAuthorFootnote{*}{Cette démonstration est due à J.P. Jouanolou.}.


